\chapter{Revisão da literatura}
Desenvolva aqui a fundamentação teórica com citações diretas e indiretas no padrão autor-data \cite{abnt10520}. Trabalhos de referência podem ser citados no corpo da frase, como \citeonline{gil2022}, ou entre parênteses ao final da oração \cite{slack2018}. Fontes eletrônicas disponíveis na internet também podem ser referenciadas normalmente no texto \cite{abntex2site}.

Para citações diretas com mais de três linhas, utilize o ambiente \texttt{citacao}, que aplica automaticamente o recuo de 4~cm da margem esquerda, fonte reduzida e espaçamento simples, dispensando aspas e exigindo a indicação da página conforme a \citeonline{abnt10520}:

\begin{citacao}
Exemplo de citação direta longa transcrita literalmente da obra consultada, redigida com extensão suficiente para ultrapassar três linhas na formatação final do documento, demonstrando o recuo de 4~cm, a fonte menor e o espaçamento simples entre linhas exigidos pelas normas vigentes \cite[p.~15]{gil2022}.
\end{citacao}
